\documentclass[10pt,a4paper]{amsart}
\usepackage[margin=19mm]{geometry}
\usepackage{amsmath,amssymb,mathtools}
\usepackage[scr=rsfs,cal=euler]{mathalfa}
\usepackage{xfrac}
\usepackage[unicode,hidelinks,bookmarks=false]{hyperref}
\newcommand{\ie}{i.e.\ }
\newcommand{\cf}{cf.\ }
\newcommand{\resp}{respectively\ }
\newcommand{\Subsection}{\subsection}
\providecommand{\nobreakdash}{\nobreak\hbox{-}}
\setlength{\emergencystretch}{2em}
\multlinegap0pt
\usepackage{amssymb}
\usepackage{mathtools}
\usepackage{bm}
\usepackage[shortlabels]{enumitem}
\usepackage{tikz}
\usepackage{xcolor}
\newtheorem{theorem}{Theorem}
\newtheorem{lemma}{Lemma}
\newcommand{\Dt}{\p_t}
\newcommand{\Lap}{\Delta}
\newcommand{\Om}{\Omega}
\newcommand{\norm}[1]{\left\lVert#1\right\rVert}
\newcommand{\ooo}{\overline}
\newcommand{\p}{\partial}
\title{Inverse source problems with reduced interior data for a coupled reaction-diffusion system}
\author{Xinyue Luo, Masahiro Yamamoto, Jin Cheng}
\date{Source: arXiv:2603.28350v1 (CC BY 4.0)}
\begin{document}
\maketitle

\noindent\emph{Source and licence.} Inverse source problems with reduced interior data for a coupled reaction-diffusion system. Version arXiv:2603.28350v1 (CC BY 4.0), distributed by arXiv under the Creative Commons Attribution 4.0 International licence, http://creativecommons.org/licenses/by/4.0/, as recorded in the arXiv licence metadata for this exact version and shown on the abstract page https://arxiv.org/abs/2603.28350v1. Full licence notice: \texttt{licenses/arxiv-2603.28350-CC-BY-4.0.txt}.\par\smallskip

\noindent\textit{Editorial scope.} Counted unit: the article's theorem thm:reference (source0.tex lines 508--526). The article's complete original proof of it is delivered with it (source0.tex lines 1017--1234, one \texttt{proof} environment of 218 lines, which the article places away from the statement). No mathematics is written, repaired or re-derived: the statement and the proof are the article's own lines, in the article's own notation, and no retained line is substituted or re-flowed.  Retained uncounted as the source-local context the proof needs: lines 101--106; lines 108--110; lines 732--738; lines 750--753; lines 797--800; lines 969--976; lines 1000--1012.  Nothing was omitted from any retained span, so the package carries no omission record.  No retained line contains a citation, so the wrapper carries no bibliography and no bibliography entry.  No retained line contains a figure environment, an image-inclusion command or drawing code, so no image is delivered and the package declares no asset dependency.  Editorial formatting only, in addition: an AMS \texttt{amsart} preamble that restates the article's own declarations and macros used by the retained lines, namely the environments \texttt{theorem}, \texttt{lemma} on separate counters, together with the standard packages those lines need.  The retained environments print in this document as Theorem 1, Lemma 1 in order of appearance, whereas the article prints the same items with the numbering of its own full document.  The complete native source of the article as distributed in the arXiv e-print for this exact version is bundled unchanged as \texttt{sources/197426/source0.tex}.
\begin{equation}\label{eq:forward}
\begin{cases}
\Dt u = d_1 \Lap u - u v^2 - a_1(x,t)\,u + f(x), & (x,t)\in Q,\\[4pt]
\Dt v = d_2 \Lap v + u v^2 - a_2(x,t)\,v, & (x,t)\in Q,
\end{cases}
\end{equation}

\begin{equation}\label{eq:bc}
u=g,\qquad v=h\qquad\text{on }\Sigma,
\end{equation}

\begin{equation}\label{eq:diffsystem}
\begin{cases}
\Dt y = d_1 \Lap y + p_1(x,t)\,y + p_2(x,t)\,z + F(x),\\[3pt]
\Dt z = d_2 \Lap z + p_3(x,t)\,z + p_4(x,t)\,y,
\end{cases}
\qquad (x,t)\in Q,
\end{equation}

\begin{equation}\label{eq:pk-bounds}
\norm{p_k}_{L^\infty(Q)}+\norm{\p_t p_k}_{L^\infty(Q)}\le M_1,
\qquad k=1,2,3,4.
\end{equation}

\begin{equation}\label{eq:alpha-max}
\alpha(x,t)\le \alpha(x,t_0)\qquad\text{for all }
(x,t)\in\Omega\times I.
\end{equation}

\begin{lemma}[Estimate on $\p_t\alpha$]\label{lem:dtalpha}
There exists $C>0$, depending on $\delta$, $\lambda$, and
$\norm{d}_{C(\overline\Om)}$, such that
\begin{equation}\label{eq:dtalpha}
|\partial_t\alpha(x,t)|\le C\,\varphi(x,t)^2
\qquad\text{for all }(x,t)\in\Omega\times I.
\end{equation}
\end{lemma}

\begin{lemma}[Carleman estimate with singular weight]\label{lem:carleman}
Let $d$ and $\alpha,\varphi$ be as above.
There exist constants $\lambda_0>0$, $s_0>0$, and $C>0$ such that for all
$\lambda\ge\lambda_0$ and $s\ge s_0$, the following holds for any
$w\in H^{2,1}(\Omega\times I)$ with $w=0$ on $\partial\Omega\times I$:
\begin{equation}\label{eq:carleman-scalar}
\mathcal I_s(w)\le
C\int_{\Omega\times I}
\bigl|\partial_t w-d_k\Delta w + B w\bigr|^2 e^{2s\alpha}\,dx\,dt
+C\int_{\omega_0\times I} s^3\varphi^3 |w|^2 e^{2s\alpha}\,dx\,dt,
\end{equation}
where $d_k\in\{d_1,d_2\}$ and $B\in L^\infty(\Om\times I)$.
\end{lemma}

\begin{theorem}[Lipschitz stability with data $(\mathcal{M}_1)$]
\label{thm:reference}
We arbitrarily choose 
a subdomain $\omega$ satisfying $\ooo{\omega} \subset\Omega$ and 
a time interval $I:= (t_1, t_2)$ with $0<t_1<t_0<t_2<T$.
Let $(u,v,f)$ and $(\widetilde u,\widetilde v,\widetilde f)$ belong to
$\mathcal U(M)$ and solve \eqref{eq:forward}--\eqref{eq:bc} with the
same coefficients and boundary data.
There exists a constant $C>0$, depending only on
$\Omega,\omega,T,t_0,t_1,t_2,M,M_0$, such that
\begin{equation}\label{eq:stab-reference}
\|f-\widetilde f\|_{L^2(\Omega)}
\le C\Bigl(
\|u(\cdot,t_0)-\widetilde u(\cdot,t_0)\|_{H^2(\Omega)}
+\sum_{\ell=0}^1 \|\partial_t^\ell(u-\widetilde u)\|_{L^2(\omega\times I)}
+\sum_{\ell=0}^1 \|\partial_t^\ell(v-\widetilde v)\|_{L^2(\omega\times I)}
\Bigr).
\end{equation}
\end{theorem}

\begin{proof}[Proof of Theorem~\ref{thm:reference}]
We divide the argument into four steps. All constants $C>0$ depend on $\Omega,\omega,T,t_0,\delta,M,M_0$ and may change from line to line.

\medskip
\noindent\textbf{Step 1: Carleman estimate for the coupled system.}

Set $y_1:=\p_t y$ and $z_1:=\p_t z$.
Differentiating \eqref{eq:diffsystem} with respect to $t$, yields
\begin{equation}\label{eq:diffsystem-dt}
\begin{cases}
\p_t y_1 = d_1\Lap y_1 + p_1 y_1 + p_2 z_1
 + (\p_t p_1) y + (\p_t p_2) z,\\[3pt]
\p_t z_1 = d_2\Lap z_1 + p_3 z_1 + p_4 y_1
 + (\p_t p_3) z + (\p_t p_4) y,
\end{cases}
\qquad (x,t)\in \Omega\times I,
\end{equation}
with $y_1=z_1=0$ on $\p\Om\times I$.

We apply Lemma~\ref{lem:carleman} to each of the four functions $y,z,y_1,z_1$:
for $y$ and $y_1$ we use diffusion constant $d_1$, and for $z$ and $z_1$ we
use $d_2$.
The right-hand sides of the Carleman estimates involve:
for $y$, $$|\p_t y-d_1\Lap y+(-p_1)y|^2=|p_2 z+F|^2;$$
for $z$, $$|\p_t z-d_2\Lap z+(-p_3)z|^2=|p_4 y|^2;$$
and similarly for $y_1$ and $z_1$ using \eqref{eq:diffsystem-dt}.

Summing the four inequalities and using \eqref{eq:pk-bounds} together with
the elementary estimate $$|A+B|^2\le2|A|^2+2|B|^2,$$ we obtain, for all
sufficiently large $s$,
\begin{align}
&\int_{\Omega\times I}\Biggl[
\frac{1}{s\varphi}\Bigl(|\p_t^2 y|^2+|\p_t^2 z|^2\Bigr)
+s^3\varphi^3\sum_{\ell=0}^1
\Bigl(|\p_t^\ell y|^2+|\p_t^\ell z|^2\Bigr)
\Biggr]e^{2s\alpha}\,dx\,dt \notag\\
&\qquad\le
C\int_{\Omega\times I}|F|^2 e^{2s\alpha}\,dx\,dt
+C\int_{\omega\times I}s^3\varphi^3\sum_{\ell=0}^1
\Bigl(|\p_t^\ell y|^2+|\p_t^\ell z|^2\Bigr)e^{2s\alpha}\,dx\,dt
\notag\\
&\qquad\quad
+C\int_{\Omega\times I}
\Bigl(|y|^2+|z|^2\Bigr)e^{2s\alpha}\,dx\,dt.
\label{eq:carleman-preabsorb}
\end{align}
The last integral arises from the lower-order terms
$(\p_t p_k)y$ and $(\p_t p_k)z$ in \eqref{eq:diffsystem-dt}.
Since the left-hand side contains
$s^3\varphi^3(|y|^2+|z|^2)e^{2s\alpha}$, choosing $s$ sufficiently large
absorbs the last integral on the right-hand side, yielding
\begin{align}
&\int_{\Omega\times I}\Biggl[
\frac{1}{s\varphi}\Bigl(|\p_t^2 y|^2+|\p_t^2 z|^2\Bigr)
+s^3\varphi^3\sum_{\ell=0}^1
\Bigl(|\p_t^\ell y|^2+|\p_t^\ell z|^2\Bigr)
\Biggr]e^{2s\alpha}\,dx\,dt \notag\\
&\qquad\le
C\int_{\Omega\times I}|F|^2 e^{2s\alpha}\,dx\,dt
+C\int_{\omega\times I}s^3\varphi^3\sum_{\ell=0}^1
\Bigl(|\p_t^\ell y|^2+|\p_t^\ell z|^2\Bigr)e^{2s\alpha}\,dx\,dt.
\label{eq:carleman-system}
\end{align}
Define the observation quantity
\begin{equation}\label{eq:D1-def}
D_1^2:=\int_{\omega\times I}\sum_{\ell=0}^1
\Bigl(|\p_t^\ell y|^2+|\p_t^\ell z|^2\Bigr)\,dx\,dt.
\end{equation}
Since $$\sup_{(x,t)\in\overline\Om\times I}
(s^3\varphi^3 e^{2s\alpha})\le e^{Cs}$$ for all $s\ge1$, we have
\begin{equation}\label{eq:obs-bound}
\int_{\omega\times I}s^3\varphi^3\sum_{\ell=0}^1
\Bigl(|\p_t^\ell y|^2+|\p_t^\ell z|^2\Bigr)e^{2s\alpha}\,dx\,dt
\le C e^{Cs} D_1^2.
\end{equation}

\medskip
\noindent\textbf{Step 2: Weighted estimate of $\p_t y(\cdot,t_0)$.}

Since $e^{2s\alpha(x,t)}\to0 $ as $ t\to t_0$, we write
\begin{align}
\int_\Om |\p_t y(x,t_0)|^2 e^{2s\alpha(x,t_0)}\,dx
&=\int_{t_0-\delta}^{t_0}\p_t\left(
\int_\Om |\p_t y|^2 e^{2s\alpha}\,dx\right)dt \notag\\
&=\int_{t_0-\delta}^{t_0}\int_\Om
\Bigl(2(\p_t y)(\p_t^2 y)+2s(\p_t\alpha)|\p_t y|^2\Bigr)
e^{2s\alpha}\,dx\,dt.
\label{eq:dt-y-identity}
\end{align}
By the Cauchy--Schwarz inequality with the factorisation
$$(\p_t y)(\p_t^2 y)=
(s\sqrt\varphi\,\p_t y)\cdot(s^{-1}\varphi^{-1/2}\,\p_t^2 y)$$
and the inequality $2ab\le a^2+b^2$, we obtain
\begin{equation}\label{eq:young-dty}
2|(\p_t y)(\p_t^2 y)|
\le s^2\varphi|\p_t y|^2+\frac{1}{s^2\varphi}|\p_t^2 y|^2.
\end{equation}
For the first term, using $\varphi\ge\delta^{-2}$, we have
$$s^2\varphi = (s^2\varphi)/(s^3\varphi^3)\cdot s^3\varphi^3
= (s\varphi^2)^{-1}\cdot s^3\varphi^3
\le C s^{-1}\cdot s^3\varphi^3.$$
For the second,
$s^{-2}\varphi^{-1}=s^{-1}\cdot(s\varphi)^{-1}$.
Thus,
\begin{equation}\label{eq:young-dty-final}
2|(\p_t y)(\p_t^2 y)|
\le \frac{C}{s}\left(s^3\varphi^3|\p_t y|^2
+\frac{1}{s\varphi}|\p_t^2 y|^2\right).
\end{equation}
For the second term in \eqref{eq:dt-y-identity}, using
Lemma~\ref{lem:dtalpha} and
$s\varphi^2\le Cs^{-1}\cdot s^3\varphi^3$ for large $s$, we obtain
$$2s|\p_t\alpha|\,|\p_t y|^2\le Cs^{-1}\cdot s^3\varphi^3|\p_t y|^2.$$

Combining with \eqref{eq:carleman-system} and \eqref{eq:obs-bound}, we
conclude
\begin{equation}\label{eq:dt-y-final}
\int_\Om |\p_t y(x,t_0)|^2 e^{2s\alpha(x,t_0)}\,dx
\le \frac{C}{s}\int_{\Om\times I}|F|^2 e^{2s\alpha}\,dx\,dt
+ C e^{Cs} D_1^2
\end{equation}
for all $s\ge s_0$.

\medskip
\noindent\textbf{Step 3: Weighted estimate of $z(\cdot,t_0)$.}

Similarly, since $e^{2s\alpha(x,t)}\to0$ as $t\to (t_0-\delta)^+$, we have
\begin{align}
\int_\Om |z(x,t_0)|^2 e^{2s\alpha(x,t_0)}\,dx
&=\int_{t_0-\delta}^{t_0}\int_\Om
\Bigl(2z\,\p_t z+2s(\p_t\alpha)|z|^2\Bigr)
e^{2s\alpha}\,dx\,dt.
\label{eq:z-identity}
\end{align}
For the first term, using the elementary inequality $2ab\le a^2+b^2$:
\[
2|z\,\p_t z|\le |z|^2+|\p_t z|^2.
\]
For the second term, Lemma~\ref{lem:dtalpha} gives
$2s|\p_t\alpha|\,|z|^2\le Cs\varphi^2|z|^2$.
Hence,
\[
2|z\,\p_t z|+2s|\p_t\alpha|\,|z|^2
\le (1+Cs\varphi^2)|z|^2+|\p_t z|^2.
\]
For $s\varphi\ge 1$ (which holds for $s$ sufficiently large, since
$\varphi\ge\delta^{-2}$), we have
\[
1+Cs\varphi^2\le \frac{C}{s}\,s^3\varphi^3,\qquad
1\le \frac{C}{s}\,s^3\varphi^3.
\]
Therefore,
\begin{equation}\label{eq:z-bound}
\int_\Om |z(x,t_0)|^2 e^{2s\alpha(x,t_0)}\,dx
\le \frac{C}{s}\int_{\Om\times I} s^3\varphi^3
\left(|z|^2+|\p_t z|^2\right)e^{2s\alpha}\,dx\,dt.
\end{equation}
The right-hand side is controlled by \eqref{eq:carleman-system}, so that
\begin{equation}\label{eq:z-final}
\int_\Om |z(x,t_0)|^2 e^{2s\alpha(x,t_0)}\,dx
\le \frac{C}{s}\int_{\Om\times I}|F|^2 e^{2s\alpha}\,dx\,dt
+ C e^{Cs} D_1^2.
\end{equation}
\medskip
\noindent\textbf{Step 4: Source identity at $t_0$ and conclusion.}

Evaluating the first equation of \eqref{eq:diffsystem} at $t=t_0$ gives
\begin{equation}\label{eq:source-identity}
F(x)=\p_t y(x,t_0)-d_1\Lap y(x,t_0)
-p_1(x,t_0)y(x,t_0)-p_2(x,t_0)z(x,t_0),
\end{equation}
whence
\begin{equation}\label{eq:F-pointwise}
|F(x)|^2\le C\Bigl(|\p_t y(x,t_0)|^2+|\Lap y(x,t_0)|^2
+|y(x,t_0)|^2+|z(x,t_0)|^2\Bigr).
\end{equation}
Define
\begin{equation}\label{eq:D2-def}
D_2:=\norm{y(\cdot,t_0)}_{H^2(\Om)}
=\norm{u(\cdot,t_0)-\tilde u(\cdot,t_0)}_{H^2(\Om)}.
\end{equation}
Multiplying \eqref{eq:F-pointwise} by $e^{2s\alpha(x,t_0)}$, integrating
over $\Om$, and using $\alpha(\cdot,t_0)$ bounded above on $\overline\Om$
for the terms involving $y(\cdot,t_0)$, we obtain
\begin{equation}\label{eq:F-weighted1}
\int_\Om |F(x)|^2 e^{2s\alpha(x,t_0)}\,dx
\le C e^{Cs} D_2^2
+C\int_\Om\Bigl(|\p_t y(x,t_0)|^2+|z(x,t_0)|^2\Bigr)
e^{2s\alpha(x,t_0)}\,dx.
\end{equation}
Substituting \eqref{eq:dt-y-final} and \eqref{eq:z-final} into
\eqref{eq:F-weighted1} yields
\begin{equation}\label{eq:F-weighted2}
\int_\Om |F(x)|^2 e^{2s\alpha(x,t_0)}\,dx
\le C e^{Cs}(D_1^2+D_2^2)
+\frac{C}{s}\int_{\Om\times I}|F|^2 e^{2s\alpha}\,dx\,dt.
\end{equation}
Since $e^{2s\alpha(x,t)}\le e^{2s\alpha(x,t_0)}$ by \eqref{eq:alpha-max}
and $F=F(x)$ is independent of $t$,
\[
\int_{\Om\times I}|F(x)|^2 e^{2s\alpha(x,t)}\,dx\,dt
\le |I|\int_\Om |F(x)|^2 e^{2s\alpha(x,t_0)}\,dx.
\]
Choosing $s\ge s_0$ so large that $C|I|/s\le 1/2$ absorbs the last term into
the left-hand side, yielding
\[
\int_\Om |F(x)|^2 e^{2s\alpha(x,t_0)}\,dx
\le C e^{Cs}(D_1^2+D_2^2).
\]
Since $\alpha(\cdot,t_0)$ is bounded below on $\overline\Om$, there exists
$c_*>0$ such that $e^{2s\alpha(x,t_0)}\ge e^{-c_* s}$ for all
$x\in\Om$.
Hence
$\norm{F}_{L^2(\Om)}^2\le C e^{Cs}(D_1^2+D_2^2)$,
and therefore
$$\norm{f-\tilde f}_{L^2(\Om)}\le C(D_1+D_2),$$
which is \eqref{eq:stab-reference}.
\end{proof}

\end{document}
